% Pursuing Stacks (cotes 134-2 to 134-8): archive cote 134-5, site page 81 (PDF page 82), typescript page 268.
% typescript. Draft compared: ps4.tex:4141-4208.
% AI-transcribed by Claude Opus 5.5 (claude-opus-5-5, adaptive thinking, effort high), Message Batches API, batch ps2 (msgbatch_01NVwr1Yuo8qJf5rykeGBXPK), 2026-11-12; not yet checked by eye.
\scanpage{134-5/81}
\mspage{268}

want. By the description of J in terms of $J_o$ , we are\inserted{now} reduced to proving
the following
\uline{Lemma}\hand{1}: Let I be an object of $M_a = M_{\underline{W}}$ . Then any two sections of I (over
$e_M$ ) are $M_c$ - homotopic.

Let, as in the previous section, $\underline{W}_o \subset \underline{W}$ be the largest of all basic
localizers satisfying Loc 4), i.e.
\[ \underline{W}_o = \hand{\{} f \in Fl((Cat)) \hand{|} \hand{\underline{\Pi}}_o(f) \text{ a bijection } \hand{\}} , \]
therefore we have
\[ M_{\underline{W}} \subset M_{\underline{W}_o} \overset{def}{=} M_o , \]
and we are reduced to proving the lemma for $\underline{W}_o$ instead of $\underline{W}$ , i.e. for $M_o$
instead of $M_{\underline{W}}$ . We'll use the\inserted{small} full subcategory C of $M_c$ generating the
contractibility structure $M_c$ , we may assume that C is stable under finite
products. Hence $C\hat{}$ is totally $\underline{W}_{oC}$ -connected i.e. totally 0-connected.
Moreover, as $C \subset M_c$ , \xout{xxxf} and $e_M$ is in C, it follows that every object
of C has a section (over $e_M = e_C$ ) - which implies that every non empty
object of $C\hat{}$ has a section, i.e. $C\hat{}$ is ``strictly totally 0-connected''
(cf. p. 144 and 149). Note that (by prop.2 b) of p. 248) we have
\[ M_o = \hand{\{} x \hand{\in} Ob\, M \hand{|} i^{\hand{*}}(x) \text{ is 0-connected in } C\hat{} \hand{\}} , \]
where i is the \xout{x}inclusion functor:
\[ i: C \hand{\longrightarrow} M . \]
We have to prove that any two section $\hand{\delta}_o,\hand{\delta}_1$ of an object I of $M_o$ are $M_c$-
homotopic, or what amounts to the same, C-homotopic. This translates
readily into the statement that $i^*(\hand{\delta}_o)$ and $i^*(\hand{\delta}_1)$ are C-homotopic in $C\hat{}$.
Thus we are reduced to the following lemma (in the case of the topos $C\hat{}$):
\uline{Lemma} 2. Let $\hand{\mathcal{C}}$\inserted{be} a totally 0-connected topos such that any non-empty object
of $\hand{\mathcal{C}}$ has a section, and let C be a small full generating subcategory,
whose elements are 0-connected. Then for any 0-connected object I of $\hand{\mathcal{C}}$ ,
and any two sections $\hand{\delta}_o,\hand{\delta}_1$ of I , these are\inserted{C-}homotopic, i.e. they can be
joined by a\inserted{finite} chain of sections, any two consecutive among which can be ob-%
tained as the images of two sections\inserted{$s_i$, $t_i$} of an \unsure{object} $x_i$ of C , by means of a
map $h_i : x_i \hand{\longrightarrow} I$ .

This lemma is essentially a restatement (cleaned from extraneous \xout{xxxxxx}
assumptions due to an awkward conceptual background) of the proposition of
page 149 (section 60), and the proof will be left to the reader.

\margin{81}
